%%%%%%%%%%%%%%%%%%%%%%%
%%% DOCUMENT SETTINGS
%%%%%%%%%%%%%%%%%%%%%%%
\documentclass[11.0pt]{exam}
\printanswers % Uncomment to show answers
\include{formatAndDefs}
%%%%%%%%%%%%%%
%%% PACKAGES
%%%%%%%%%%%%%%
\usepackage{amsmath,amsfonts,amssymb,amsthm} % math fonts and symbols
\usepackage{bbm} % Math blackboard font
\usepackage{enumitem} % For reference enumerations lists
\usepackage{textcomp} % some symbols (like copyright)
\usepackage[top = 2cm, bottom = 2cm, right = 1.2cm, left = 1.2cm, paper = a4paper, headheight = 6cm]{geometry} % Margins setup
\usepackage{xcolor} % Colors
%%%%%%%%%%%%%%%%%%%%%%%%%%
%%% MACROS AND COMMANDS
%%%%%%%%%%%%%%%%%%%%%%%%%%
\definecolor{CUNEForange}{RGB}{237, 114, 1}
\newcommand{\N}{\mathbb{N}} % natural number set
\newcommand{\R}{\mathbb{R}} % real number set
%%%%%%%%%%%%%%%%%%%%%%%%%%%%%%%%%
%%% HEADER AND FOOT 
%%%%%%%%%%%%%%%%%%%%%%%%%%%%%%%%%
% Defining information
\newcommand{\degree}[1]{Bachelor in Philosophy, Politics and Economics}
\newcommand{\shortdeg}[1]{PPE}
\newcommand{\exam}[1]{1st Midterm}
\newcommand{\subject}[1]{Mathematics}
\newcommand{\theyear}[1]{2023/24}
\newcommand{\ccauthor}[1]{A. Azze \textcopyright\the\year}
\newcommand{\thedate}[1]{Oct 18, 2023}
\newcommand{\university}[1]{CUNEF University}
% Setting header
\usepackage{fancyhdr}
\pagestyle{fancy}
\renewcommand{\headrulewidth}{0.1mm}
\renewcommand{\footrulewidth}{0.1mm}
% \renewcommand{\sectionmark}[1] % To avoid sections' titles being displayed on the header
{\markright{\MakeUppercase{}}}
\setlength\headsep{0.5cm}
% Header and footer
\fancyhead[L]{\degree{} \\ \subject{} - \theyear{}}
\fancyhead[R]{\thedate{} \\ \ccauthor{}}
\fancyfoot[L]{\university{}}
\fancyfoot[C]{}
\fancyfoot[R]{\thepage}

%%%%%%%%%%%%%%%%%%%%%%%%%
%%% DOCUMENT CONTENT
%%%%%%%%%%%%%%%%%%%%%%%%%
\begin{document}
%--- Instructions ---%
\begin{center}

    {\Huge \textbf{\exam{}}}
    \vspace{0.7cm}

    \begin{flushright}
        \textbf{Time Limit:} 50 minutes
    \end{flushright}
    \vspace{-0.2cm}
    
    \fbox{\parbox{\textwidth}{\centering 
        \begin{enumerate}[leftmargin = 0.5cm, label = $\bullet$, rightmargin = 0.2cm]
    
            \item \textbf{Every non-trivial calculation and claim in the exam must be properly justified.}

            \item The use of calculators is not necessary and therefore not allowed.

            \item Digital devices such as mobile phones, tablets, and laptops, as well as internet access, are forbidden.

            \item You can use physical resources such as books, notes, and other printed materials to assist you during the exam.
 
        \end{enumerate}
    }}
\end{center}
%--- Instructions ---%

%--- Questions ---%

\begin{questions}

% Question
\question[3]
    Consider the function $f(x) = \frac{18x^3 - 3x^2 + 5x + 1}{6x^3 - 24x}$.
    \begin{enumerate}[leftmargin = 0.5cm, label = (\alph*), rightmargin = 0.2cm]
        \item Calculate the limit $\lim_{x\rightarrow\infty} f(x)$.
        \item For $g(x) = f(x) + \ln(x)/\sqrt{x}$ obtain the limit $\lim_{x\rightarrow\infty} g(x)$.
        \item For $h(x)$ such that $f(x) \leq h(x) \leq g(x)$, calculate the limit $\lim_{x\rightarrow\infty} h(x)$.
    \end{enumerate}
% Solution
\begin{solution}[0cm]
    \begin{enumerate}[leftmargin = 0.5cm, label = (\alph*), rightmargin = 0.2cm]
        \item Dividing by $x^3$ both in the numerator and denominator of $f(x)$, we get that
        \begin{align*}
            \lim_{x\rightarrow\infty} f(x) = \frac{18x^3 - 3x^2 + 5x + 1}{6x^3 - 24x} = \lim_{x\rightarrow\infty} f(x) = \frac{18 - 3x^{-1} + 5x^{-2} + 1x^{-3}}{6 - 24x^{-2}} = 2
        \end{align*}
        \item By using L'Hôspital's rule, we obtain 
        \begin{align*}
            \lim_{x\rightarrow\infty} \frac{\ln(x)}{\sqrt{x}} = \lim_{x\rightarrow\infty} \frac{2\sqrt{x}}{x} = 0.
        \end{align*}
        Hence, the linearity of the limit operator implies that $\lim_{x\rightarrow\infty} g(x) = \lim_{x\rightarrow\infty}f(x) = 2$.
        \item Since $f(x) \leq h(x) \leq g(x)$, then 
        \begin{align*}
            2 = \lim_{x\rightarrow\infty}f(x) \leq \lim_{x\rightarrow\infty}h(x) \leq \lim_{x\rightarrow\infty}g(x) = 2.
        \end{align*}
        Hence, $\lim_{x\rightarrow\infty}h(x) = 2$.
    \end{enumerate}
\end{solution}

% Question
\question[2]
    Compute the derivative of the function $f(x) = x + (x + x^2)^2$ at $x = 1$ and use it to propose an approximation of $f(1.01)$.
%% Solution
\begin{solution}[0cm]
    By applying the linearity of the differential operator and the chain rule, we get that 
    \begin{align*}
        f'(x) = 1 + 2(x + x^2)(x + x^2)' = 1 + 2(x + x^2)(1 + 2x).
    \end{align*}
    Due to the definition of the derivative, we can propose the approximation
    \begin{align*}
        f(1.01) \approx f'(1)(1 - 1.01) + f(1),
    \end{align*}
    Where $f'(1) = 1 + 2\times2\times 3 = 13$ and $f(1) = 5$. Then,
    \begin{align*}
        f(1.01) \approx 13\times0.01 + 5 = 0.13 + 5 = 5.13.
    \end{align*}
    The actual value of $f(1.01)$, rounded to 4 decimal places, is $5.1313$, so you can appreciate the accuracy of the approximation based on the first derivative.
\end{solution}

% Question
\question[1]
    Obtain the product rule of differentiation by using the quotient rule.
%% Solution
\begin{solution}[0cm]
    Let $f(x)$ and $g(x)$ be two functions. Assume that $g(x)\neq 0$ for all $x$. Then
    \begin{align*}
        (f(x)g(x))' = \left(\frac{f(x)}{1/g(x)}\right)' &= \frac{f'(x)/g(x) + f(x)g'(x)/g^2(x)}{1/g^2(x)} \\
        &= g^2(x)(f'(x)/g(x) + f(x)g'(x)/g^2(x)) = f'(x)g(x) + f(x)g'(x).
    \end{align*}
\end{solution}

% Question
\question[1]
    Let $f(x)$ be a convex function. Is $g(x) = e^{f(x)}$ convex?
%% Solution
\begin{solution}[0cm]
    By applying the chain rule and the product rule, we get
    \begin{align*}
        g'(x) = f'(x)e^{f(x)}, \quad g''(x) = f''(x)e^{f(x)} + (f'(x))^2e^{f(x)}.
    \end{align*}
    Therefore, $g''(x) \geq 0$, meaning that $g(x)$ is convex.
\end{solution}

% % Question
% \question[2]
% The tobacco market is a monopoly, with supply and demand sets
% \begin{align*}
%     S = \{(q, p) : 15p − 2q = 20\},\quad D = \{(q, p) : q + 5p = 40\}.
% \end{align*}
% Given that the government has imposed an excise tax of $T$ dollars to the tabaco market:

% \begin{enumerate}[leftmargin = 0.5cm, label = (\alph*), rightmargin = 0.2cm]
        
%     \item Define the governments revenue function in terms of the excise tax and the number of units sold.

%     \item Find the excise tax that maximizes the government revenue.
        
% \end{enumerate}
% % Solution
% \begin{solution}[0cm]
%     Solution
% \end{solution}

% % Question
% \question[2]
%     A firm manufactures two goods, $X$ and $Y$. The market price of these goods is unaffected by the level of the firm's production. The firm's cost function is $C(x, y) = 2x^2+ xy + 2y^2$. The market prices for $X$ and $Y$ are $12$ and $18$ dollars per unit, respectively. Determine the number of units of each good that the firm should produce to maximize its profit.  
% % Solution
% \begin{solution}[0cm]
%     Solution
% \end{solution}

% % Question
% \question[2]
%     A wine dealer has certain amount of units of barrels of wine, which he can sell today at $K$ dollars or wait for a better deal. The price of his barrels of wine increases in time according the the function $Kt^n$, for $n \in \N$. Assume that the market holds an interest rate of $r \R_+$ on continuously compounding levels, that is, $x$ today's dollars equal $xe^{-r t}$ dollars $t$ units of time ahead.
    
%     \begin{enumerate}[leftmargin = 0.5cm, label = (\alph*), rightmargin = 0.2cm]

%         \item Define the profit made by selling the barrels of wine at time $t$.
    
%         \item When should the dealer sell the barrels of wine to maximize his profit? Hint: you can use Problem 3 to ease the optimization.
    
%     \end{enumerate}
% % Solution
% \begin{solution}[0cm]
%     Solution
% \end{solution}

% Question
\question[3]
    A firm's production function is defined as $P(k, l) = k^2l$, where $k$ stands for capital invested and $l$ for labour input. At time $t\in\R_+$, the capital invested and labour input are given by
    $$
    k(t) = 4 + 3t, \quad l(t) = 9e^{-t}.
    $$

    \begin{enumerate}[leftmargin = 0.5cm, label = (\alph*), rightmargin = 0.2cm]
        
        \item How does a small change in time affect production?

        \item When would the production be maximized?

        % \item Define the log-production function as $f(t) = \log(P(k(t), l(t)))$. When $f(t)$ reaches its maxima?
        
        % \item Proof that the production is also maximized at the times where $f(t)$ reaches its maxima.

        % \item What is the maximum production level the firm can achieve?
        
    \end{enumerate}

% Solution
\begin{solution}[0cm]
    Let $f(t) = P(k(t), l(t))$. By using the chain rule for two variables, and the fact that $k'(t) = 3$ and $l'(t) = -9e^{-t}$, and $P_k(k, l) = 2kl$ and $P_l(k, l) = k^2$, we get
    \begin{align*}
        f'(t) &= k'(t)P_k(k(t), l(t)) + l'(t)P_l(k(t), l(t)) \\
        &= 6k(t)l(t) - 9e^{-t}k^2(t) \\
        &= 6(4 + 3t)9e^{-t} - 9e^{-t}(4 + 3t)^2 \\
        &= 9e^{-t}(4 + 3t)(2 - 3t).
    \end{align*}
    Therefore, for $t \geq 0$, $f'(t) = 0$ if and only if $2 - 3t = 0$, meaning that $t = 2/3$. It is easy to check that $f'(t) > 0$ for every $t\in[0, 2/3)$, and $f'(t) < 0$ for every $t \in (2/3, \infty)$. Hence, $f(t)$ reaches its global maximum at $t = 2/3$.
\end{solution}

\end{questions}	

\end{document}